\documentclass[12pt]{article}

%% ---- Packages ----
\usepackage{amsmath, amsthm, amssymb}
\usepackage{mathtools}
\usepackage{tikz-cd}
\usepackage[T1]{fontenc}
\usepackage[utf8]{inputenc}
\usepackage{microtype}
\usepackage{enumitem}
\usepackage{listings}
\usepackage{xcolor}
\usepackage{geometry}
\usepackage{natbib}
\usepackage{hyperref}  % load last among packages
\geometry{margin=1in}

%% ---- Lean listings style ----
\lstdefinelanguage{Lean4}{
  keywords={def,theorem,lemma,structure,where,by,let,have,show,exact,refine,
            intro,intros,apply,rw,simp,congr,ext,cases,rcases,induction,
            open,universe,variable,instance,class,namespace,end,import,
            fun,match,if,then,else,return,do,sorry,calc,suffices,obtain,
            use,constructor,fconstructor,fapply,dsimp,erw,conv,at,from},
  sensitive=true,
  comment=[l]{--},
  morecomment=[s]{/-}{-/},
  morestring=[b]",
  morekeywords=[2]{Type,Prop,Sort,true,false},
}
\lstset{
  language=Lean4,
  basicstyle=\ttfamily\small,
  keywordstyle=\color{blue!70!black}\bfseries,
  keywordstyle=[2]\color{teal},
  commentstyle=\color{gray}\itshape,
  stringstyle=\color{brown},
  breaklines=true,
  frame=single,
  backgroundcolor=\color{gray!6},
  numbers=left,
  numberstyle=\tiny\color{gray},
  numbersep=5pt,
}

%% ---- Theorem environments ----
\newtheorem{theorem}{Theorem}[section]
\newtheorem{lemma}[theorem]{Lemma}
\newtheorem{proposition}[theorem]{Proposition}
\newtheorem{corollary}[theorem]{Corollary}
\theoremstyle{definition}
\newtheorem{definition}[theorem]{Definition}
\newtheorem{example}[theorem]{Example}
\newtheorem{remark}[theorem]{Remark}
\newtheorem{notation}[theorem]{Notation}

%% ---- Math macros ----
\newcommand{\Cat}{\mathbf{Cat}}
\newcommand{\op}{\mathrm{op}}
\newcommand{\id}{\mathrm{id}}
\newcommand{\Id}{\mathrm{Id}}
\newcommand{\StrongTrans}{\mathbf{StrongTrans}}
\newcommand{\yoneda}{\mathcal{Y}}
\newcommand{\yo}{\yoneda_0}
\newcommand{\Hom}{\mathrm{Hom}}
\newcommand{\obj}{\mathrm{obj}}
\newcommand{\map}{\mathrm{map}}
\newcommand{\B}{\mathcal{B}}
\newcommand{\bicatprod}{\times}
\newcommand{\psfun}{\Rightarrow_{\mathrm{ps}}}


%% ---- Title info ----
\title{A Lean 4 Formalization of the Bicategorical Yoneda Lemma}
\author{Spencer Woolfson}
\date{\today}

\begin{document}

\maketitle

\begin{abstract}
  We present a formalization of the Yoneda lemma for bicategories in Lean 4, using
  the \texttt{Mathlib} library.  The classical Yoneda lemma for ordinary categories
  states a natural bijection $\Hom(y(b), F) \cong F(b)$ between natural transformations
  out of the representable functor $y(b)$ and elements of $F(b)$.  In the bicategorical
  setting this bijection is upgraded to an equivalence of categories, naturally in both
  the representing object $b \in \B^{\op}$ and the pseudofunctor
  $F : \B^{\op} \to \Cat$.  Concretely, we construct pseudofunctors
  $\mathtt{yonedaPairing}$ and $\mathtt{yonedaEvaluation}$ and exhibit them as
  equivalent via an explicit pair of strong transformations together with unit and
  counit isomorphisms.  A nontrivial aspect of the formalization is universe management:
  the two sides of the equivalence naturally land in different universe levels, requiring
  an auxiliary universe-lifting pseudofunctor $\mathtt{CatPseudoULift}$.
\end{abstract}

\tableofcontents
\newpage

%%============================================================
\section{Introduction}
%%============================================================

The Yoneda lemma is one of the most fundamental results in category theory.  In its
classical form it states that, for a locally small category $\mathcal{C}$ and a functor
$F : \mathcal{C}^{\op} \to \mathbf{Set}$, there is a natural bijection
\[
  \Hom_{\mathbf{Set}^{\mathcal{C}^{\op}}}(y(c), F) \;\cong\; F(c)
\]
where $y : \mathcal{C} \to \mathbf{Set}^{\mathcal{C}^{\op}}$ is the Yoneda embedding.
The lemma has wide-ranging consequences: it identifies objects with the functors they
represent, it underlies the theory of adjunctions, and it is the engine behind most
universal constructions in algebra and geometry.

In the bicategorical setting, the appropriate notion of ``functor'' is a
\emph{pseudofunctor} (or homomorphism of bicategories), natural transformations are
replaced by \emph{strong (pseudo-natural) transformations}, and the Yoneda bijection is
promoted to an \emph{equivalence of categories}: for a bicategory $\B$ with Yoneda
embedding $\yoneda : \B \to [\B^{\op}, \Cat]_{\mathrm{ps}}$, we have
\begin{equation}\label{eq:bi-yoneda}
  \StrongTrans(\yo(b),\, F) \;\simeq\; F(b)
\end{equation}
naturally in $b \in \B^{\op}$ and $F : \B^{\op} \to \Cat$.

This paper reports on a machine-checked proof of~\eqref{eq:bi-yoneda} carried out in
Lean~4 using Mathlib~\citep{mathlib4}.  The formalization closely follows the internal
categorical description found in~\citet{johnson-yau2021} and~\citet{leinster2004}:
rather than constructing the equivalence object-by-object, we exhibit it as a
biequivalence $\mathtt{yonedaPairing} \simeq \mathtt{yonedaEvaluation}$ of
pseudofunctors $\B^{\op} \times [\B^{\op}, \Cat] \to \Cat$
(Section~\ref{sec:main}).

\medskip\noindent\textbf{Organization.}
Section~\ref{sec:background} recalls the relevant background on bicategories,
pseudofunctors, and strong transformations, fixing notation.
Section~\ref{sec:yoneda-embedding} reviews the bicategorical Yoneda embedding as it
appears in Mathlib.
Section~\ref{sec:main} states and proves the main theorem.
Section~\ref{sec:universes} discusses the universe issues that arise in the
formalization and the auxiliary construction $\mathtt{CatPseudoULift}$.
Section~\ref{sec:lean} highlights design choices in the Lean code.
Section~\ref{sec:open} discusses directions for future formalization work.

%%============================================================
\section{Background}\label{sec:background}
%%============================================================

\subsection{Bicategories}

We work with \emph{bicategories} in the sense of B{\'e}nabou~\citep{benabou1967}.

\begin{definition}
  A \emph{bicategory} $\B$ consists of:
  \begin{enumerate}[label=(\roman*)]
    \item a collection of \emph{objects} $\obj(\B)$;
    \item for each pair of objects $a, b \in \B$, a category $\B(a,b)$ whose objects
          are \emph{1-morphisms} and whose morphisms are \emph{2-morphisms};
    \item \emph{horizontal composition} functors
          $\circ : \B(b,c) \times \B(a,b) \to \B(a,c)$;
    \item \emph{identity} 1-morphisms $\id_a \in \B(a,a)$;
    \item natural isomorphisms (the \emph{associator}
          $\alpha_{f,g,h} : (h \circ g) \circ f \cong h \circ (g \circ f)$
          and \emph{unitors} $\lambda_f : \id \circ f \cong f$,
          $\rho_f : f \circ \id \cong f$) satisfying the pentagon and triangle axioms.
  \end{enumerate}
\end{definition}

In Lean/Mathlib the bicategory axioms are bundled in the \lstinline{Bicategory} typeclass.
We use the notation $f \triangleright \sigma$ (\lstinline{whiskerLeft}) and
$\sigma \triangleleft f$ (\lstinline{whiskerRight}) for horizontal composition of a
1-morphism with a 2-morphism.

\subsection{Pseudofunctors}

\begin{definition}
  A \emph{pseudofunctor} $F : \B \to \mathcal{C}$ between bicategories consists of:
  \begin{enumerate}[label=(\roman*)]
    \item an object map $F_0 : \obj(\B) \to \obj(\mathcal{C})$;
    \item for each $a,b \in \B$, a functor $F_{a,b} : \B(a,b) \to \mathcal{C}(Fa, Fb)$;
    \item coherence isomorphisms $F(\id_a) \cong \id_{Fa}$ and
          $F(g \circ f) \cong F(g) \circ F(f)$, subject to coherence axioms.
  \end{enumerate}
\end{definition}

In Mathlib pseudofunctors are called \lstinline{Pseudofunctor} and are written
\lstinline{B ⥤ᵖ C} in Lean.

\subsection{Strong Transformations}

\begin{definition}
  Let $F, G : \B \to \mathcal{C}$ be pseudofunctors.  A \emph{strong (pseudo-natural)
  transformation} $\eta : F \to G$ consists of:
  \begin{enumerate}[label=(\roman*)]
    \item for each object $a \in \B$, a 1-morphism $\eta_a : Fa \to Ga$;
    \item for each 1-morphism $f : a \to b$ in $\B$, an invertible 2-morphism
          (the \emph{naturality isomorphism})
          $\eta_f : G(f) \circ \eta_a \cong \eta_b \circ F(f)$;
    \item coherence conditions with respect to identities and composition.
  \end{enumerate}
\end{definition}

The category whose objects are strong transformations $F \to G$ and whose morphisms are
\emph{modifications} is denoted $\StrongTrans(F, G)$ and written
\lstinline{Pseudofunctor.StrongTrans F G} in Lean.

%%============================================================
\section{The Bicategorical Yoneda Embedding}\label{sec:yoneda-embedding}
%%============================================================

Fix a bicategory $\B$.  The \emph{Yoneda embedding} for bicategories is a pseudofunctor
\[
  \yoneda : \B \;\longrightarrow\; [\B^{\op}, \Cat]_{\mathrm{ps}}
\]
where $[\B^{\op}, \Cat]_{\mathrm{ps}}$ denotes the bicategory of pseudofunctors
$\B^{\op} \to \Cat$.  This is provided by Mathlib as \lstinline{Bicategory.yoneda} (and
the component at an object $b$ is written \lstinline{yoneda₀ b}).

Concretely, $\yo(b) : \B^{\op} \to \Cat$ is the pseudofunctor that
\begin{itemize}
  \item sends an object $a \in \B^{\op}$ to the hom-category $\B(a, b)$;
  \item sends a 1-morphism $f : a' \to a$ in $\B^{\op}$ (i.e.\ $f : a \to a'$ in $\B$)
        to the precomposition functor $f^* : \B(a, b) \to \B(a', b)$,
        $g \mapsto g \circ f$.
\end{itemize}

%%============================================================
\section{The Main Theorem}\label{sec:main}
%%============================================================

\subsection{The Two Pseudofunctors}

\begin{definition}[{\lstinline{yonedaPairing}}]\label{def:pairing}
  The \emph{pairing pseudofunctor}
  \[
    P : \B^{\op} \times [\B^{\op}, \Cat]_{\mathrm{ps}} \;\longrightarrow\; \Cat
  \]
  is defined by:
  \begin{align*}
    P(b, F) &= \StrongTrans(\yo(b), F), \\
    P(f : b' \to b,\; \alpha : F \to G)(\eta) &= \alpha \circ \eta \circ \yo(f).
  \end{align*}
\end{definition}

\begin{definition}[{\lstinline{yonedaEvaluation}}]\label{def:evaluation}
  The \emph{evaluation pseudofunctor}
  \[
    E : \B^{\op} \times [\B^{\op}, \Cat]_{\mathrm{ps}} \;\longrightarrow\; \Cat
  \]
  is defined by:
  \begin{align*}
    E(b, F) &= F(b), \\
    E(f : b' \to b,\; \alpha : F \to G)(s) &= G(f)(\alpha_b(s)).
  \end{align*}
\end{definition}

\subsection{The Equivalence}

The main result is:

\begin{theorem}[Bicategorical Yoneda Lemma]\label{thm:yoneda}
  There is a biequivalence
  \[
    P \;\simeq\; E
  \]
  in the bicategory of pseudofunctors
  $\B^{\op} \times [\B^{\op}, \Cat]_{\mathrm{ps}} \to \Cat$,
  i.e.\ strong transformations
  $\Phi : P \to E$ and $\Psi : E \to P$ together with invertible modifications
  $\Phi \circ \Psi \cong \Id_E$ and $\Psi \circ \Phi \cong \Id_P$.
\end{theorem}

The equivalence is witnessed by the following data.

\begin{definition}[{\lstinline{yonedaLemmaForwards}}]
  The strong transformation $\Phi : P \to E$ has component
  \[
    \Phi_{(b,F)} : \StrongTrans(\yo(b), F) \;\longrightarrow\; F(b), \qquad
    \eta \;\longmapsto\; \eta_b(\id_b).
  \]
  That is, we evaluate the strong transformation $\eta$ at the object $b$ and then
  apply the resulting functor to the identity morphism $\id_b \in \B(b,b)$.
\end{definition}

\begin{definition}[{\lstinline{yonedaLemmaBackwards}}]
  The strong transformation $\Psi : E \to P$ has component
  \[
    \Psi_{(b,F)} : F(b) \;\longrightarrow\; \StrongTrans(\yo(b), F), \qquad
    s \;\longmapsto\; \bigl[(a, f) \mapsto F(f)(s)\bigr].
  \]
  That is, the strong transformation assigned to $s \in F(b)$ sends an object
  $a \in \B^{\op}$ and a morphism $f \in \B(a, b)$ to $F(f)(s)$.
\end{definition}

\begin{proposition}
  The composites $\Phi \circ \Psi$ and $\Psi \circ \Phi$ are isomorphic to the
  respective identities via the modifications $\mathtt{yonedaInvHomId}$ and
  $\mathtt{yonedaHomInvId}$.  These are assembled from:
  \begin{itemize}
    \item The \emph{unit}: for $\eta \in \StrongTrans(\yo(b), F)$, the component
          isomorphism at $f : a \to b$ is the composite
          $(\eta.\text{naturality}(f))^{-1}|_{\id_b} \circ \eta_a.\map(\rho_f)$,
          where $\rho_f : f \circ \id_b \cong f$ is the right unitor.
    \item The \emph{counit}: for $s \in F(b)$, the component is
          $F.\text{mapId}(b)|_s$, the identity coherence of $F$ applied to $s$.
  \end{itemize}
\end{proposition}

%%============================================================
\section{Universe Considerations}\label{sec:universes}
%%============================================================

A subtle issue arises when comparing $P$ and $E$ as pseudofunctors into $\Cat$:
they naturally land in different universe levels.

Fix universe variables $u$, $v$, $w$ so that $\B$ has objects in $\mathtt{Type}\ u$,
hom-categories with morphisms in $\mathtt{Type}\ v$, and 2-morphisms in
$\mathtt{Type}\ w$.  Then:
\begin{itemize}
  \item $E'(b, F) = F(b)$ lands in $\Cat^{w,v}$ (i.e.\ a category with
        hom-sets in $\mathtt{Type}\ w$ and objects in $\mathtt{Type}\ v$).
  \item $P(b, F) = \StrongTrans(\yo(b), F)$ lands in
        $\Cat^{\max(u,v,w), \max(u,v,w)}$.
\end{itemize}
These two universe levels do not match, so the naive statement
$P \simeq E'$ is not type-correct.

To resolve this we define an auxiliary pseudofunctor
\[
  \mathtt{CatPseudoULift} : \Cat^{v_1, u_1} \;\to\; \Cat^{\max(v_1,v_2),\, \max(u_1,u_2)}
\]
which promotes every small category to a larger universe by applying
$\mathtt{ULift}$ on objects and $\mathtt{ULiftHom}$ on hom-sets.  The actual evaluation
pseudofunctor used in the theorem is
\[
  \mathtt{yonedaEvaluation} = \mathtt{yonedaEvaluation'} \;;\; \mathtt{CatPseudoULift},
\]
where $\mathtt{yonedaEvaluation'}$ is the raw evaluation functor landing in
$\Cat^{w,v}$.  Universe-lifting is lossless: the equivalence
$\mathtt{CatLiftIso}$ witnesses $C \simeq \mathtt{CatLift}(C)$ for every category $C$.

%%============================================================
\section{Design Choices in the Lean Formalization}\label{sec:lean}
%%============================================================

\subsection{The \texttt{BiEquiv} Structure}

Rather than working with adjoint equivalences (which carry triangle identity data), we
introduce a lighter structure sufficient for stating the Yoneda result:

\begin{lstlisting}
structure BiEquiv (x y : B) where
  map : x ⟶ y
  inv : y ⟶ x
  homInvId : map ≫ inv ≅ 𝟙 x
  invHomId : inv ≫ map ≅ 𝟙 y
\end{lstlisting}

This records an internal equivalence in a bicategory --- a pair of 1-morphisms with
invertible 2-morphisms witnessing the round-trips --- without requiring the triangle
identities of an adjoint equivalence.  Every adjoint equivalence gives a
\lstinline{BiEquiv}, but the converse requires additional work.

\subsection{Opacity and Definitional Reductions}

Several auxiliary lemmas use \lstinline{erw} (extended rewrite) rather than \lstinline{rw}
because the relevant equalities hold only after definitional unfolding of universe-lift
wrappers such as \lstinline{ULiftHom.up} and \lstinline{ULift.upFunctor}.  Replacing
these with dedicated simp lemmas that expose the required reductions is a natural
cleanup.

\subsection{Modifications and Extensionality}

The extensionality lemma

\begin{lstlisting}
lemma StrongTrans.Hom.ext {m1 m2 : η ⟶ θ} (h : m1.as = m2.as) : m1 = m2
\end{lstlisting}

is used throughout to reduce goals about 2-morphisms in $\StrongTrans(F, G)$ to
equalities of the underlying modifications.  This avoids unfolding the single-field
wrapper \lstinline{StrongTrans.Hom} by hand and mirrors the standard
\lstinline{Functor.ext} pattern already present in Mathlib.

%%============================================================
\section{Future Work}\label{sec:open}
%%============================================================

The formalization established here opens several directions for further development
of higher category theory in Lean~4 and Mathlib.

\paragraph{Biadjunctions.}
The biequivalence of Theorem~\ref{thm:yoneda} is the core ingredient in showing that
the Yoneda embedding $\yoneda : \B \to [\B^{\op}, \Cat]_{\mathrm{ps}}$ is a
\emph{biadjunction} --- the appropriate bicategorical analogue of a fully faithful
functor admitting a left adjoint.  Formalizing the general theory of biadjunctions
in Mathlib, including the bicategorical adjoint functor theorem, would be a natural
continuation~\citep{lack2010}.

\paragraph{The Grothendieck construction for bicategories.}
The classical Grothendieck construction builds a fibration over a category $\mathcal{C}$
from a functor $\mathcal{C} \to \Cat$.  For bicategories there is an analogous
\emph{bicategorical Grothendieck construction}, producing a bicategration over $\B$ from
a pseudofunctor $\B \to \Cat$~\citep{street1974}.  The pseudofunctors
$\mathtt{yonedaPairing}$ and $\mathtt{yonedaEvaluation}$ are precisely the inputs this
construction takes, so a formalization of bicategorical fibrations would sit naturally on
top of the present work.

\paragraph{Pseudomonads and pseudoalgebras.}
A \emph{pseudomonad} on a bicategory $\B$ is a monoid object in the hom-bicategory of
pseudoendofunctors, with composition and unit isomorphisms satisfying coherence
conditions.  Pseudomonads model computational effects at a bicategorical level and
include important examples such as the free cocompletion monad on $\Cat$.  The Yoneda
lemma is the key tool for showing that the free cocompletion of $\B$ is the bicategory
of presheaves $[\B^{\op}, \Cat]_{\mathrm{ps}}$.

\paragraph{Tricategories and coherence.}
Bicategories, pseudofunctors, strong transformations, and modifications themselves
assemble into a \emph{tricategory}~\citep{gordon-power-street1995}.  Formalizing
tricategories in Lean and proving the Gordon--Power--Street coherence theorem --- every
tricategory is triequivalent to a Gray-category --- is a significant open problem in
proof assistant libraries and would extend the hierarchy already present in Mathlib.

\paragraph{Profunctors and enriched structures.}
The bicategory $\mathbf{Prof}$ of small categories and profunctors (distributors) is one
of the central objects of study in formal category theory.  The Yoneda lemma implies
that the Yoneda embedding $y : \mathcal{C} \to [\mathcal{C}^{\op}, \mathbf{Set}]$ is the
unit of the free cocompletion monad on $\mathbf{Prof}$, and the induced theory of
\emph{Cauchy completion} and \emph{absolute colimits} is a natural target for
formalization~\citep{fiore-gambino-hyland-winskel2008}.

\paragraph{Yoneda structures.}
\citet{street1974} introduced \emph{Yoneda structures} as an axiomatic framework that
isolates the properties of the Yoneda embedding needed to develop formal category theory
in an arbitrary 2-category.  Formalizing Yoneda structures in Lean would unify the
treatment of ordinary, enriched, and internal category theory under a single abstract
framework, allowing theorems to be proved once and instantiated to many settings.

%%============================================================
\section*{Acknowledgements}
%%============================================================

% Add acknowledgements here.

\bibliographystyle{plainnat}
\bibliography{references}

\end{document}
